\section{Supuestos}\label{an:supuestos}
\begin{enumerate}
\item \textbf{Solo automóviles.} No hay buses, motos ni tranvía. El tranvía hacía abortar SUMO~1.12 en los cruces, y los buses no cambian la velocidad media que se predice.
\item \textbf{Comportamiento igual.} Los vehículos conectados conducen igual que los demás (modelo Krauss): estar conectado no implica conducción automatizada.
\item \textbf{Rutas fijas.} Sin el aviso, los conductores no se desvían, por recomendación del docente. En el SIL, el cambio de ruta es precisamente la intervención, y solo lo hacen los vehículos avisados.
\item \textbf{Penetración V2X fija.} Es de 50~\% en los días normales y de 30 y 70~\% en los de estrés. Los vehículos conectados son una muestra aleatoria del tráfico.
\item \textbf{Comunicación ideal.} El DENM llega sin pérdidas ni retardo a todos los vehículos V2X del área de relevancia; la cobertura de radio no se simula.
\item \textbf{Incidentes simulados.} Un incidente es un vehículo detenido por carril bloqueado. Su inicio y fin reales se toman de la salida de SUMO.
\item \textbf{Lluvia.} Se representa con más demanda y conductores más prudentes, no con un modelo de adherencia.
\item \textbf{Días independientes.} Cada día empieza con la red vacía y una hora de calentamiento que se descarta.
\end{enumerate}

\section{Parámetros del generador sintético y su fuente}\label{an:param}
La Tabla~\ref{tab:param} resume los parámetros del generador. Donde no hay una cifra publicada para Cuenca, el valor se fijó por un criterio propio que se indica (\emph{diseño}).

\begin{table}[h]
\centering
\caption{Parámetros del generador sintético}
\label{tab:param}
\scriptsize
\setlength{\tabcolsep}{2.5pt}
\begin{tabular}{@{}L{2.2cm}L{2.25cm}L{3.95cm}@{}}
\toprule
Parámetro & Valor & Fuente o criterio \\
\midrule
Red vial & OSM, 15 parroquias urbanas & OpenStreetMap; límites de OSM o Ley de Tránsito (50~km/h urbano); Américas 50/70~km/h~\cite{canafe2026irap} \\
Composición del parque & solo automóviles & 32~\% de los viajes son en vehículo privado~\cite{cazorla2026unimobility}; buses y motos fuera del alcance (Anexo~\ref{an:supuestos}) \\
Demanda pico & V0 = 17\,000 veh/h & Diseño: punto de saturación de la red (eficiencia ≥ 0,95), barrido con 3 semillas \\
Viajes por día & \textasciitilde200\,000 & Conservador frente a \textasciitilde400\,000 viajes en auto estimados (1,5~M viajes, 32~\% en auto, 1,2 personas por auto~\cite{elmercurio2026movilidad,cazorla2026unimobility,elmercurio2020carros}) \\
Perfil horario & picos 07--08, 12--14, 18--19 & Picos medidos por radar~\cite{jacome2024radar} y pico de mediodía~\cite{pena2021canaris} \\
Factor del día & 0,88--1,07 & Diseño: variabilidad entre días sin colapsar la red \\
Conducción & Krauss, σ 0,5, τ 1,0 s, speedFactor N(1; 0,1) & Valores por defecto de SUMO~\cite{sumo2026vtypes} \\
Lluvia & 20~\% de días; velocidad deseada ×0,88; demanda +4~\% & Lluvia fuerte: −6 a −9~\% de velocidad libre y −10 a −11~\% de capacidad~\cite{hranac2006weather} \\
Duración de incidentes & 10--25 min & Incidentes ``menores'' (< 30 min)~\cite{fhwa2009mutcd6i,li2018incidentduration} \\
Penetración V2X & 50~\% (30 y 70 en estrés) & Diseño: escenario intermedio y sensibilidad \\
Semillas & 1000 + día; 2026; 2027 & Reproducibilidad \\
\bottomrule
\end{tabular}
\end{table}


\section{Validación del dato sintético frente a datos reales}\label{an:valida}
Como no hay datos de la EMOV EP para el periodo, la validación compara magnitudes del dataset con cifras publicadas de Cuenca (Tabla~\ref{tab:valida}; \texttt{validar\_sintetico.py}). El volumen simulado de la Av.~de las Américas se mide como un aforo: los vehículos que cruzan una línea transversal en cada tramo de 200~m, sumando ambos sentidos, la vía central y las laterales, en los 24 laborables secos.

\begin{table}[h]
\centering
\caption{Dato sintético frente a cifras reales de Cuenca}
\label{tab:valida}
\scriptsize
\setlength{\tabcolsep}{2.5pt}
\begin{tabular}{@{}L{2.4cm}L{2.0cm}L{2.15cm}L{1.75cm}@{}}
\toprule
Magnitud & Real & Sintético \sint & Lectura \\
\midrule
Américas, hora pico & 4500--5400 veh/h~\cite{elmercurio2021americas,elmercurio2022americas} & máx. 2610 veh/h (p90 2291) & \textasciitilde50--58~\% \\
Américas, por día & 75\,000 veh~\cite{elmercurio2021americas} & máx. 34\,831 veh & \textasciitilde46~\% \\
Viajes en auto por día & \textasciitilde400\,000 (estimado) & 207\,784 & \textasciitilde52~\% \\
Hora pico mañana & 07:00--08:00~\cite{jacome2024radar} & 08:00--09:00 & 1~h de desfase \\
Hora pico tarde & 18:00--19:00~\cite{jacome2024radar} & 18:00--19:00 & coincide \\
Efecto de la lluvia & −6 a −9~\% vel. libre, −10 a −11~\% capacidad~\cite{hranac2006weather} & −14,8~\% vel. de viaje & más severo \\
Límites, Américas & 50 y 70 km/h~\cite{canafe2026irap} & 50 y 70 km/h (y otros tramos) & coincide \\
\bottomrule
\end{tabular}
\end{table}

\textbf{Lectura.} El sintético reproduce la forma del tráfico real: el pico de la tarde y, con una hora de desfase, el de la mañana; los límites de velocidad del corredor; y un efecto de la lluvia del mismo signo, algo más severo porque combina menor velocidad y más demanda. En volumen está en el orden de magnitud correcto, pero es conservador: los tres indicadores de volumen quedan en torno a la mitad de las cifras reales, que además son de prensa y sin método declarado. La demanda pico se fijó en la saturación de la red simulada sin rerouting; con más demanda la congestión no se disuelve. Por eso los resultados deben leerse como sintéticos: la mejora relativa frente a la línea base es la cifra relevante, y con demanda real el beneficio del aviso podría ser distinto. Un aforo propio en la Av. de las Américas, cargado con el formato del Anexo~\ref{an:transicion}, permitiría recalibrar la demanda.


\section{Transición al dato real}\label{an:transicion}
El código no depende del origen de los datos, sino de un formato fijo de tablas Parquet, documentado en \texttt{FORMATO\_DATOS.md}:
\begin{itemize}
\item velocidades por celda cada 5~min;
\item conteos, velocidad y ocupación por arista de corredor cada minuto;
\item calendario de días;
\item registro de incidentes.
\end{itemize}
El script \texttt{validar\_datos.py} revisa que una carpeta tenga esas tablas, columnas, tipos y rangos físicos, y la variable \texttt{RVH\_DATOS} apunta el proyecto a esa carpeta. Cuando lleguen los conteos de la EMOV EP o de una cámara, basta con convertirlos a ese formato. Los datos de una sola fuente se copian en las dos vistas (\texttt{gt} y \texttt{obs}). Ningún script de entrenamiento, evaluación o visualización cambia.
